\documentclass[11pt]{article}
\usepackage[T1]{fontenc}
\usepackage{lmodern}
\usepackage[margin=1in]{geometry}
\usepackage{microtype}
\usepackage{xcolor}
\usepackage{amsmath,amssymb}
\usepackage[hidelinks]{hyperref}
\setlength{\parindent}{0pt}
\setlength{\parskip}{0.7em}
\definecolor{accent}{HTML}{234E70}
\definecolor{panel}{HTML}{EDF4F8}
\newcommand{\statusbox}[1]{%
  \begin{center}
  \colorbox{panel}{\parbox{0.91\linewidth}{\centering\color{accent}\bfseries #1}}
  \end{center}}

\begin{document}

{\Large\bfseries Literature answer: Gross' Brownian motion\\[-0.1em]
fails the polarity principle\par}

\vspace{0.35em}
{\color{accent}\rule{\linewidth}{1.1pt}}

\statusbox{Status: literature\_already\_answered (negative answer)}

\section*{Original open problem}

L. Beznea, A. Cornea, and M. R\"ockner, \emph{Potential theory of infinite
dimensional L\'evy processes}, arXiv:1007.2379v2, Section 4.8, page 19,
immediately after Remark 4.15, ask:

\begin{quote}
Does the axiom of polarity hold for infinite-dimensional Brownian motion?
\end{quote}

Here the axiom is Hunt's Hypothesis (H): every semipolar Borel set is polar.

\section*{Decisive supporting paper}

P. J. Fitzsimmons, \emph{Gross' Brownian Motion Fails to Satisfy the Polarity
Principle}, Rev. Roumaine Math. Pures Appl. \textbf{59} (2014), no. 1,
87--91. Theorem 1, on journal page 87 (PDF page 1), states that Gross'
Brownian motion does not satisfy the axiom of polarity. Proposition 1, on
journal page 88, supplies a Borel set that is semipolar but not polar.

\section*{Identification}

The supporting paper explicitly identifies the same question. Its opening
paragraph cites the Beznea--Cornea--R\"ockner paper, says that it left open
whether Gross' process satisfies the axiom of polarity, and announces a
counterexample.

The mechanism is an intrinsic clock. On the classical Wiener state space
$E=C([0,1];\mathbb R)$, Fitzsimmons defines the quadratic-variation functional
\[
 Q(x)=\lim_{n\to\infty}\sum_{k=1}^{2^n}
 \bigl(x(k2^{-n})-x((k-1)2^{-n})\bigr)^2
\]
when the limit exists. The shell $S=\{Q=1\}$ is not polar because a process
started at the origin belongs to $S$ at time one almost surely. Proposition 2
proves
\[
 Q(X_t)=Q(X_0)+t \qquad (t>0),
\]
whenever $Q(X_0)<\infty$. Consequently the process visits $S$ at most once;
the strong Markov argument in the paper makes $S$ semipolar. Thus $S$ is the
required semipolar, nonpolar counterexample.

\section*{Scope and provenance}

The formal theorem is stated for Gross' Brownian motion on the classical
Wiener space. This already disproves general validity of the polarity axiom
for infinite-dimensional Brownian motion. Fitzsimmons also states that the
construction can be modified for a general abstract Wiener space, but does
not spell out that extension. The paper does not claim the analogous result
for the broader class of infinite-dimensional L\'evy processes.

This is an explicit later-paper answer, not a new result of the present run.
The identification was confirmed from the original page-19 question, the
supporting paper's Theorem 1 and Propositions 1--2, an exact-phrase search for
``axiom of polarity'' with ``infinite dimensional Brownian motion'', and the
2019 survey on Hunt's Hypothesis (H).

\begin{thebibliography}{9}
\bibitem{bcr}
L. Beznea, A. Cornea, and M. R\"ockner,
\emph{Potential theory of infinite dimensional L\'evy processes},
J. Funct. Anal. \textbf{261} (2011), 2845--2876;
arXiv:1007.2379.

\bibitem{fitzsimmons}
P. J. Fitzsimmons,
\emph{Gross' Brownian Motion Fails to Satisfy the Polarity Principle},
Rev. Roumaine Math. Pures Appl. \textbf{59} (2014), no. 1, 87--91.

\bibitem{survey}
Z.-C. Hu and W. Sun,
\emph{Hunt's Hypothesis (H) for Markov Processes: Survey and Beyond},
arXiv:1908.06825.
\end{thebibliography}

\end{document}
